\documentclass[11pt]{article}

% ---------- packages ----------
\usepackage[utf8]{inputenc}
\usepackage[T1]{fontenc}
\usepackage{lmodern}
\usepackage[margin=1in]{geometry}
\usepackage{amsmath,amssymb}
\usepackage{booktabs}
\usepackage{array}
\usepackage{enumitem}
\usepackage{xcolor}
\usepackage{tcolorbox}
\usepackage{longtable}
\usepackage[round]{natbib}
\usepackage[hidelinks]{hyperref}

\newtcolorbox{keybox}[1]{colback=blue!4,colframe=blue!45!black,
  fonttitle=\bfseries,title=#1,boxrule=0.6pt,left=4pt,right=4pt,top=3pt,bottom=3pt}

% ---------- notation shortcuts ----------
\newcommand{\gF}{g_F}
\newcommand{\gB}{g_B}
\newcommand{\hF}{h_F}
\newcommand{\hB}{h_B}
\newcommand{\fF}{f_F}
\newcommand{\fB}{f_B}
\newcommand{\Cstar}{C^{*}}
\newcommand{\GMX}{G_{\mathrm{MX}}}
\newcommand{\lb}{\mathit{lb}}

\title{\vspace{-1.5em}Bidirectional Heuristic Search on Voxel Maps\\[0.3em]
\large Background from the Literature and a Project Plan}
\author{Search Methods in Artificial Intelligence --- Ben-Gurion University\\
Advisor: Lior Siag}
\date{\today}

\begin{document}
\maketitle
\vspace{-1em}

\begin{abstract}
\noindent This document distills the six theory papers and one domain paper assigned by
Lior Siag into exactly what is needed to start the project: run bidirectional heuristic
search algorithms (A$^{*}$, MM, BAE$^{*}$, BiA$^{*}$, GBFS, plus unidirectional variants)
on the Warthog 3D voxel-map benchmarks inside the HOG2 framework, and measure how close
each algorithm comes to the theoretical minimum number of node expansions. It is organized
so the theory builds up in the order Lior recommended --- the ``must-expand'' foundation
first, then the lower-bound machinery, then the concrete algorithms --- and closes with the
voxel domain details that drive the loader, the experimental methodology, and a recommended
two-person workflow.
\end{abstract}

\tableofcontents

% =====================================================================
\section{Project in one page}
\label{sec:overview}

We are searching for a least-cost path of cost $\Cstar$ between a \texttt{start} and a
\texttt{goal} state. \emph{Bidirectional heuristic search} (BiHS) runs two searches at
once: a \textbf{forward} search from \texttt{start} and a \textbf{backward} search from
\texttt{goal} (using reversed moves). The promise is that two searches, each travelling
``half the distance,'' can expand far fewer nodes than a single unidirectional
$A^{*}$~\citep{holte2017mm}.

The concrete goal (per Lior's brief) is to \textbf{connect the Warthog voxel-map benchmark
to the bidirectional-search code in HOG2}, then run every required algorithm from both
directions and analyze the results. HOG2 already contains the algorithms
(\texttt{MM.h}, \texttt{BAE.h}, \texttt{TemplateAStar.h}, \texttt{NBS.h},
\texttt{BidirectionalGreedyBestFirst.h}) and a voxel domain (\texttt{Voxels.\{h,cpp\}}),
so the real work is (a) a loader from the Warthog voxel format into HOG2's domain, and
(b) an experiment driver that runs all algorithms both ways and records metrics.

Everything in the assigned papers answers two questions:
\begin{enumerate}[leftmargin=1.4em,itemsep=1pt]
  \item \textbf{Theory:} which nodes must \emph{any} correct algorithm expand? This gives a
        provable lower bound --- the yardstick we compare against
        (\S\ref{sec:mustexpand}--\S\ref{sec:bounds}).
  \item \textbf{Algorithms:} how do MM, BAE$^{*}$, A$^{*}$, BiA$^{*}$ and GBFS try to reach
        that minimum, and when does each win? (\S\ref{sec:algorithms}).
\end{enumerate}

% =====================================================================
\section{Background and notation}
\label{sec:notation}

Learn this notation first --- it appears in every paper.
\begin{itemize}[leftmargin=1.4em,itemsep=2pt]
  \item $\gF(n)$: cost of the best known path \texttt{start}$\to n$ (forward).
        $\gB(n)$: cost of the best known path $n\to$\texttt{goal} (backward).
  \item $\hF(n)$: heuristic estimate of $n\to$\texttt{goal}.
        $\hB(n)$: estimate of \texttt{start}$\to n$.
  \item $\fF(n)=\gF(n)+\hF(n)$ and $\fB(n)=\gB(n)+\hB(n)$, exactly as in $A^{*}$.
  \item $\Cstar$: optimal solution cost. $U$: cost of the best (incumbent) solution found so
        far. $\varepsilon$: the cheapest edge cost in the graph.
\end{itemize}

\begin{keybox}{Front-to-end vs.\ front-to-front (know the difference)}
All algorithms in this project are \textbf{front-to-end}: each side's heuristic points at
the \emph{fixed opposite endpoint} ($\hF$ aims at \texttt{goal}, $\hB$ aims at
\texttt{start}). This contrasts with \textbf{front-to-front} heuristics that estimate the
distance between two \emph{arbitrary} frontier nodes --- more informative but rarely
practical~\citep{eckerle2017sufficient}. The whole must-expand theory below is stated for
the front-to-end setting.
\end{keybox}

\paragraph{Admissible vs.\ consistent.}
A heuristic is \emph{admissible} if it never overestimates ($\hF(n)\le d(n,\texttt{goal})$).
It is \emph{consistent} if it additionally satisfies the triangle inequality
$\hF(n)\le d(n,n')+\hF(n')$. \textbf{Our voxel heuristics (3D-octile / Euclidean distance)
are consistent}, not merely admissible --- and consistency unlocks tighter lower bounds
(\S\ref{sec:consistency}), so it is the regime we actually live in.

% =====================================================================
\section{The must-expand theory: which nodes are unavoidable}
\label{sec:mustexpand}

This is the intellectual spine, developed across
\citet{eckerle2017sufficient}, \citet{chen2017nbs}, and the two Shaham
papers~\citep{shaham2017minimal,shaham2018consistent}. Understand it and the rest follows.

\subsection{The must-expand \emph{pair} \citep{eckerle2017sufficient}}
In unidirectional $A^{*}$ the obligation is on \emph{single nodes}: $A^{*}$ must expand
every node with $f<\Cstar$~\citep{dechter1985generalized}. Bidirectional search is
different --- the obligation is on \textbf{pairs} $(u,v)$, where $u$ is a forward-frontier
node and $v$ a backward-frontier node.

Define the lower bound on the cost of any path forced through \emph{both} $u$ and $v$:
\begin{equation}
  \lb(u,v)=\max\bigl\{\,\fF(u),\ \fB(v),\ \gF(u)+\gB(v)\,\bigr\}.
  \label{eq:lb}
\end{equation}
The three terms are three independent reasons such a path must cost at least this much:
$\fF(u)$ is the forward estimate of a solution through $u$; $\fB(v)$ is the backward
estimate through $v$; and $\gF(u)+\gB(v)$ is the \emph{meeting term} --- the guaranteed cost
of reaching $u$ from the start plus reaching the goal from $v$, \emph{even if the connecting
piece were free}. The meeting term is the uniquely bidirectional ingredient.

\begin{keybox}{The rule}
$(u,v)$ is a \textbf{must-expand pair} iff $\lb(u,v)<\Cstar$. The obligation is discharged
by expanding \emph{either} $u$ forward \emph{or} $v$ backward --- not necessarily both.
\end{keybox}

\emph{Why (intuition).} An expansion-based, black-box algorithm cannot ``see'' edges it has
not generated. If it skips both $u$ and $v$, an adversary can insert one cheap edge $u\to v$
that creates a solution cheaper than $\Cstar$; the algorithm never notices and returns a
sub-optimal answer. So to be \emph{guaranteed} optimal it must touch at least one side of
every must-expand pair. \citet{eckerle2017sufficient} also show this generalizes the classic
$A^{*}$ result~\citep{dechter1985generalized}, which is the special case $\hB=0$.

\subsection{The must-expand graph and minimum vertex cover \citep{chen2017nbs}}
Turn this into a bipartite graph $\GMX$: forward nodes on the left, backward nodes on the
right, an edge between $u$ and $v$ iff $(u,v)$ is a must-expand pair. Because expanding a
node ``covers'' every edge touching it, and every must-expand edge must be covered:
\begin{keybox}{The theoretical floor}
The minimum number of expansions any BiHS algorithm can achieve equals the
\textbf{Minimum Vertex Cover (MVC) of $\GMX$}. Since $\GMX$ is bipartite, the MVC is
computable in polynomial time (K\H{o}nig's theorem / max-flow), so this floor is
\emph{measurable per instance} --- it is the denominator we report every algorithm against.
\end{keybox}
\citet{chen2017nbs} also introduced \textbf{NBS} (Near-optimal Bidirectional Search), which
provably expands at most $2\times$MVC and is worst-case optimal --- our natural
near-optimal baseline.

\subsection{Hitting the floor: fractional MM \citep{shaham2017minimal}}
\citet{shaham2017minimal} show the MVC need not be built explicitly: cluster nodes by
$g$-value, and the optimal cover is always a clean split --- all lowest-$g$ forward nodes up
to a cutoff, plus all lowest-$g$ backward nodes up to a dual cutoff --- computable in time
linear in the number of distinct $g$-values. They introduce \textbf{fMM($p$)}
(``fractional MM''), which meets at fraction $p$ of the path: $p{=}1$ is forward $A^{*}$,
$p{=}0$ is reverse $A^{*}$, $p{=}\tfrac12$ is MM. There is an optimal $p^{*}$ for which
fMM($p^{*}$) expands \emph{exactly} the MVC --- but $p^{*}$ depends on $\Cstar$ and the
instance, so it is only knowable \emph{after} solving. This is the permanent gap between the
theoretical optimum and any real algorithm.

\begin{keybox}{Experimental consequence}
The best meeting point is domain-dependent and \emph{usually not the middle}. Do not assume
MM's $\Cstar/2$ split is optimal on voxel maps --- measure the distribution of $p^{*}$.
\end{keybox}

\subsection{What consistency buys \citep{shaham2018consistent}}
\label{sec:consistency}
Because our heuristics are \emph{consistent}, the heuristic difference itself becomes a free
front-to-front lower bound on the ``gap'' $d(u,v)$ between the forward node $u$ and the
backward node $v$. Define
\begin{equation}
  h_C(u,v)=\max\bigl\{\,\hF(u)-\hF(v),\ \hB(v)-\hB(u),\ \varepsilon\,\bigr\},
  \label{eq:hc}
\end{equation}
and the must-expand condition tightens to $\gF(u)+\gB(v)+h_C(u,v)<\Cstar$. The two heuristic
terms are directional mirror images that each lower-bound the \emph{same} forward gap
$d(u,v)$: consistency of $\hF$ gives $\hF(u)-\hF(v)\le d(u,v)$, and consistency of $\hB$
along the same $u\to v$ segment gives $\hB(v)-\hB(u)\le d(u,v)$. (On an \emph{undirected}
graph such as a voxel map, where $d(u,v)=d(v,u)$, both signs hold and the terms may be
replaced by absolute values $|\hF(u)-\hF(v)|$, $|\hB(u)-\hB(v)|$ --- the ``$\mathit{lb}_{CU}$''
variant; on directed graphs only the signed form is admissible.) The extra non-negative term
makes the inequality fail more often, so \emph{fewer pairs are must-expand}
$\Rightarrow$ sparser $\GMX$ $\Rightarrow$ smaller floor. In plain terms, consistency lets
you \emph{prove} more pairs are safe to skip, which is precisely what consistency-exploiting
algorithms (BAE$^{*}$) capitalize on. \citet{shaham2018consistent}
prove a matching optimal algorithm exists (per heuristic-equivalence class), confirming the
lower bound is tight.

% =====================================================================
\section{Unifying the lower bounds \citep{siag2025bridging}}
\label{sec:bounds}

\citet{siag2025bridging} --- Lior's own capstone paper, and the recommended starting read
--- proves that all the practical ``lower bound'' formulas invented over the years are just
different subsets of the single $\max\{\cdot\}$ expression in Eq.~\eqref{eq:lb}, and that
infinitely many valid bounds can be generated from it. Two derived per-node quantities are
central:
\begin{align}
  d_D(n) &= g_D(n)-h_{\bar D}(n) && \text{(\emph{heuristic error}: how much the opposite
                                     heuristic underestimated),}\\
  b_D(n) &= f_D(n)+d_D(n)        && \text{(the \emph{$b$-value}, BAE$^{*}$'s priority).}
\end{align}
where $D$ is a direction and $\bar D$ its opposite. From these come the family of global
bounds, all valid lower bounds on $\Cstar$; the most useful in practice are:
\begin{itemize}[leftmargin=1.4em,itemsep=1pt]
  \item the \textbf{g-bound} $g\mathrm{Min}_F+g\mathrm{Min}_B+\varepsilon$;
  \item the \textbf{Kaindl--Kainz bounds}~\citep{kaindl1997reconsidered}
        $f\mathrm{Min}_F+d\mathrm{Min}_B$ and $f\mathrm{Min}_B+d\mathrm{Min}_F$;
  \item the \textbf{$b$-bound} $\tfrac12\bigl(b\mathrm{Min}_F+b\mathrm{Min}_B\bigr)$,
        the one BAE$^{*}$ pushes.
\end{itemize}
No single bound dominates the others, but the empirical headline is decisive for us:

\begin{keybox}{Practical guidance from Siag \& Shperberg (2025)}
\begin{enumerate}[leftmargin=1.3em,itemsep=1pt]
  \item \textbf{Default to BAE$^{*}$} under consistent heuristics --- best balance of node
        expansions and runtime.
  \item Existing algorithms already come within $\sim$5--15\% of the MVC floor, so there is
        little room left to reduce expansions by inventing new bounds.
  \item \textbf{On grid/map-like domains the consistency benefit is small}, and reverse
        $A^{*}$ sometimes beats everything because of directional asymmetry. Whether the
        fancy bidirectional algorithm wins on voxel maps is an \emph{empirical} question ---
        which is exactly what this project answers.
\end{enumerate}
\end{keybox}

% =====================================================================
\section{The algorithms we will run}
\label{sec:algorithms}

\begin{center}
\renewcommand{\arraystretch}{1.3}
\begin{longtable}{@{}p{2.1cm}p{2.2cm}p{5.4cm}p{1.4cm}@{}}
\toprule
\textbf{Algorithm} & \textbf{Type} & \textbf{Priority / key idea} & \textbf{Optimal?}\\
\midrule
$A^{*}$ (+ reverse) & unidirectional & expand min $f=g+h$ & yes\\
MM & bidirectional & expand min $\mathit{pr}=\max(f,2g)$ (meet in the middle) & yes\\
BiA$^{*}$ & bidirectional & bidirectional $A^{*}$, min $f$ each side; frontiers not capped & yes\\
BAE$^{*}$ & bidirectional & expand min $b=f+d$, $d=g-h_{\text{opp}}$; exploits heuristic error & yes\\
GBFS & greedy & expand min $h$ & \textbf{no}\\
NBS & bidirectional & near-optimal; $\le 2\times$MVC (baseline) & yes\\
\bottomrule
\end{longtable}
\end{center}

\subsection{MM --- guaranteed to meet in the middle \citep{holte2017mm}}
The whole trick is the priority function
\begin{equation}
  \mathit{pr}_F(n)=\max\bigl(\fF(n),\ 2\,\gF(n)\bigr)
  \qquad(\text{backward symmetric}).
  \label{eq:mm}
\end{equation}
A node is expanded only when its priority $\le\Cstar$, so the $2g$ term forces
$\gF(n)\le\Cstar/2$: \emph{neither frontier ever ventures past the halfway distance from its
own root}. (This is about how far each search reaches, not physically where they meet.)
Implementation essentials:
\begin{itemize}[leftmargin=1.4em,itemsep=1pt]
  \item two full search states (Open/Closed/$g$ per direction); backward uses reversed moves;
  \item each step expands whichever direction owns the global minimum priority
        $C=\min(\mathit{pr}\text{min}_F,\mathit{pr}\text{min}_B)$;
  \item \textbf{solution detection}: when a generated node already sits on the opposite Open
        list, a full path exists --- update $U\leftarrow\min(U,\gF(n)+\gB(n))$;
  \item \textbf{termination}: stop when
        $U\le\max\!\bigl(C,\ f\mathrm{min}_F,\ f\mathrm{min}_B,\
        g\mathrm{min}_F+g\mathrm{min}_B+\varepsilon\bigr)$;
        \emph{not} at first frontier contact (that yields sub-optimal answers);
  \item \textbf{MM0} ($h{=}0$, expand min $2g$) is the bidirectional brute-force baseline,
        analogous to Dijkstra.
\end{itemize}
MM wins with moderate-strength heuristics; a very strong heuristic favors $A^{*}$, a very
weak one favors MM0.

\subsection{BAE$^{*}$ --- bidirectional $A^{*}$ with error \citep{alcazar2020unifying,siag2025bridging}}
BAE$^{*}$ is a bidirectional best-first search whose priority is the $b$-value
$b_D(n)=f_D(n)+d_D(n)=f_D(n)+\bigl(g_D(n)-h_{\bar D}(n)\bigr)$, terminating on the $b$-bound.
Intuitively it ``charges the heuristic's looseness back into'' a node's priority, delaying
loose nodes. Under consistency $b$ stays monotone, so BAE$^{*}$ is optimal. The
implementation delta from BiA$^{*}$ is small (swap the priority, use the $b$-bound to stop),
and it is Lior's recommended default. \citet{alcazar2020unifying} showed it can expand up to
an order of magnitude fewer ``necessarily expanded'' nodes than near-optimal algorithms such
as NBS on some domains, because exploiting heuristic inaccuracy turns many would-be
must-expand nodes into delayable ones.

\subsection{GBFS --- greedy best-first}
GBFS expands by lowest $h$ only and is \textbf{not cost-optimal}; none of the
must-expand/MVC theory applies to it. Report it on a separate axis (fast, but its path cost
and expansions are not comparable against the MVC floor).

% =====================================================================
\section{The voxel domain \citep{nobes2023voxel}}
\label{sec:voxel}

This is the practical layer the loader must match.

\paragraph{What a voxel map is.} A 3D integer lattice $\text{width}\times\text{height}\times
\text{depth}$; each voxel is filled (obstacle) or free, with up to 26 neighbours (the 3D
analogue of the 8-neighbour 2D grid). The benchmark ships three families: \textbf{Sandstone}
(dense porous rock, all $400^3$, ${\sim}80\%$ filled), \textbf{Descent} (1995 game levels,
rooms + corridors, up to ${\sim}1.3$ billion voxels), and \textbf{Industrial Plants}
(pipe-routing layouts, sparse). In total: 46 maps, 92{,}000 instances.

\begin{keybox}{File format --- the single most error-prone step}
\begin{verbatim}
voxel                    <- OR "rev_voxel"
<width> <height> <depth>
x0 y0 z0
x1 y1 z1
...
\end{verbatim}
\textbf{\texttt{voxel}} $\Rightarrow$ listed coordinates are the \emph{filled/obstacle}
voxels (everything else free). \textbf{\texttt{rev\_voxel}} $\Rightarrow$ listed coordinates
are the \emph{free} voxels (everything else blocked) --- used for dense maps like Sandstone.
\emph{Branch on the first token, or dense maps come out completely inverted.} The domain
paper names the fields but gives no byte-exact spec: verify indexing order, 0-\ vs
1-based, whitespace, and scenario columns against a real downloaded file and the Warthog
loader source before trusting a parser.
\end{keybox}

\paragraph{Movement model (match exactly for comparable costs).}
\begin{itemize}[leftmargin=1.4em,itemsep=1pt]
  \item \textbf{26-connected} moves with 3D-octile costs: straight (face) $=1$,
        edge-diagonal $=\sqrt2$, corner-diagonal $=\sqrt3$. Hence $\varepsilon=1$.
  \item \textbf{No corner-cutting}: a diagonal move is legal only if the voxels it clips are
        all free.
  \item Region \emph{reachability} is computed with 6-connectivity (flood fill); keep that
        distinct from the 26-connected \emph{move} model.
  \item ``Without diagonals'' (per the brief) $=$ the 6 face moves at unit cost; ``with
        diagonals'' $=$ full 26-connected. Since the domain supports diagonals, we run both.
\end{itemize}

\paragraph{Scenario files and metrics.} Each map ships $\sim$2{,}000 \texttt{start}/\texttt{goal}
voxel triples, stored in the minimum-cost direction (which controls directional bias ---
directly relevant when validating bidirectional symmetry), likely carrying an optimal cost.
The currency of comparison is \textbf{node expansions}. Report two flavors used by the
papers: expansions to \emph{prove} optimality (bound $<\Cstar$) versus total ($\le\Cstar$),
each as a ratio to the per-instance MVC floor. Also log runtime and solution cost --- all
optimal algorithms must return the same cost as $A^{*}$, which is our correctness check.

% =====================================================================
\section{What to expect on our maps}
\label{sec:expect}

Synthesizing the empirical findings across the papers, on grid/voxel-like domains:
\begin{itemize}[leftmargin=1.4em,itemsep=1pt]
  \item $A^{*}$ (and reverse $A^{*}$) are surprisingly competitive; strong Euclidean/octile
        heuristics narrow the bidirectional advantage~\citep{holte2017mm,siag2025bridging}.
  \item MM/BAE$^{*}$ pay off most with moderate heuristics; the consistency benefit is
        modest on maps~\citep{shaham2018consistent,siag2025bridging}.
  \item The best meeting point is rarely the exact middle; directional asymmetry matters,
        so always run \emph{both} directions~\citep{shaham2017minimal,nobes2023voxel}.
  \item GBFS is fast but sub-optimal --- separate axis.
  \item Report every optimal algorithm as a ratio to the MVC floor (or to NBS's
        $\le2\times$MVC line) rather than as raw counts.
\end{itemize}

% =====================================================================
\section{Recommended workflow and two-person split}
\label{sec:workflow}

The work factors cleanly into an \textbf{infrastructure track} (get maps into HOG2 and run
things) and an \textbf{algorithms/analysis track} (make all algorithms run both ways and
turn numbers into findings). We propose \textbf{Person~A = Infrastructure} and
\textbf{Person~B = Algorithms \& Analysis}, with shared milestones and paired reviews at
each phase boundary.

\subsection{Phased plan}
\renewcommand{\arraystretch}{1.25}
\begin{longtable}{@{}p{0.6cm}p{3.2cm}p{5.0cm}p{2.6cm}@{}}
\toprule
\textbf{Ph} & \textbf{Milestone} & \textbf{Deliverable / exit criterion} & \textbf{Owner}\\
\midrule
0 & Setup \& read & HOG2 built; benchmarks fetched; both read papers 2, 3, 1, 6 & Both\\
1 & Format spike & One real Sandstone (\texttt{rev\_voxel}) + one Descent (\texttt{voxel}) map parsed and visualized; format documented & A (lead), B (review)\\
2 & Loader & Warthog $\to$ HOG2 \texttt{Voxels} loader; 26-conn + no-corner-cutting + octile costs; unit tests & A\\
3 & Baseline run & $A^{*}$ and reverse $A^{*}$ produce optimal costs matching benchmark; metrics logged (expansions, runtime, cost) & B (lead), A (loader support)\\
4 & Algorithm sweep & MM, MM0, BiA$^{*}$, BAE$^{*}$, GBFS, NBS wired into the driver; run both directions, with and without diagonals & B\\
5 & MVC yardstick & Per-instance MVC of $\GMX$ computed; expansion ratios produced & B (lead), A (perf/IO)\\
6 & Scale \& robustness & Large Descent/Plant maps stream without OOM; timeouts/memory caps handled & A\\
7 & Analysis \& report & Tables + plots per family; crossover discussion; final write-up & Both\\
\bottomrule
\end{longtable}

\subsection{Responsibility split}
\begin{center}
\renewcommand{\arraystretch}{1.25}
\begin{longtable}{@{}p{5.7cm}p{5.7cm}@{}}
\toprule
\textbf{Person A --- Infrastructure} & \textbf{Person B --- Algorithms \& Analysis}\\
\midrule
Warthog voxel loader (\texttt{voxel}/\texttt{rev\_voxel}), scenario parser & Experiment driver looping all algorithms $\times$ both directions $\times$ diagonal on/off\\
HOG2 \texttt{Voxels} domain wiring: 26-connectivity, no-corner-cutting, octile costs & Configure/verify MM, MM0, BiA$^{*}$, BAE$^{*}$, GBFS, NBS in HOG2\\
Bit-packed storage \& streaming for billion-voxel maps; build system & Metrics: expansions ($<\Cstar$ and $\le\Cstar$), runtime, cost, meeting point\\
Correctness \& perf harness; timeouts, memory caps, reproducible runs & MVC-of-$\GMX$ computation as the theoretical floor; ratio reporting\\
CI-style regression on a small map subset & Statistical analysis, plots, and the results tables\\
\bottomrule
\end{longtable}
\end{center}

\subsection{Ways of working}
\begin{itemize}[leftmargin=1.4em,itemsep=1pt]
  \item \textbf{Contract early.} Agree the loader's output data structure (the occupancy grid
        + domain interface) in Phase~1 so A and B can work in parallel behind it. B can start
        on a tiny hand-made map before the real loader lands.
  \item \textbf{Correctness gate.} No algorithm result is trusted until its returned cost
        equals $A^{*}$'s on the same instance (Phase~3 establishes the oracle).
  \item \textbf{Reference implementation.} Mirror the SPL-BGU \texttt{BiHS-Direction-Choosing}
        usage example and the \texttt{BiHS-Consistent-F2E} repo for how HOG2's BiHS code is
        driven.
  \item \textbf{Cross-review at phase boundaries.} Each phase's exit criterion is reviewed by
        the other person; this catches the two classic bugs early --- inverted
        \texttt{rev\_voxel} polarity, and corner-cutting sneaking in cheaper-than-optimal
        paths.
  \item \textbf{Scale last.} Develop on Sandstone ($400^3$) and small Descent maps; only move
        to billion-voxel maps once the pipeline is correct and instrumented.
\end{itemize}

% =====================================================================
\section*{Reading order (Lior's recommendation)}
\addcontentsline{toc}{section}{Reading order}
Start with \citet{siag2025bridging} (the unifying big picture); if clear, continue to
\citet{eckerle2017sufficient} (the must-expand foundation). Then
\citet{holte2017mm} (MM --- the most concrete algorithm), \citet{alcazar2020unifying}
(BAE$^{*}$, the likely default), and \citet{shaham2017minimal,shaham2018consistent} for the
minimal-set and consistency depth. Read \citet{nobes2023voxel} alongside the loader work.

% =====================================================================
\begin{thebibliography}{10}

\bibitem[Alc\'azar et~al.(2020)]{alcazar2020unifying}
Alc\'azar, V., Riddle, P., and Barley, M. (2020).
\newblock A Unifying View on Individual Bounds and Heuristic Inaccuracies in
  Bidirectional Search.
\newblock In \emph{Proc.\ AAAI Conference on Artificial Intelligence}, pp.\ 2327--2334.

\bibitem[Chen et~al.(2017)]{chen2017nbs}
Chen, J., Holte, R.~C., Zilles, S., and Sturtevant, N.~R. (2017).
\newblock Front-to-End Bidirectional Heuristic Search with Near-Optimal Node
  Expansions.
\newblock In \emph{Proc.\ IJCAI}, pp.\ 489--495.

\bibitem[Dechter and Pearl(1985)]{dechter1985generalized}
Dechter, R. and Pearl, J. (1985).
\newblock Generalized best-first search strategies and the optimality of A$^{*}$.
\newblock \emph{Journal of the ACM}, 32(3):505--536.

\bibitem[Eckerle et~al.(2017)]{eckerle2017sufficient}
Eckerle, J., Chen, J., Sturtevant, N.~R., Zilles, S., and Holte, R.~C. (2017).
\newblock Sufficient Conditions for Node Expansion in Bidirectional Heuristic
  Search.
\newblock In \emph{Proc.\ ICAPS}.

\bibitem[Holte et~al.(2017)]{holte2017mm}
Holte, R.~C., Felner, A., Sharon, G., and Sturtevant, N.~R. (2017).
\newblock MM: A bidirectional search algorithm that is guaranteed to meet in
  the middle.
\newblock \emph{Artificial Intelligence}, 252:232--266.

\bibitem[Kaindl and Kainz(1997)]{kaindl1997reconsidered}
Kaindl, H. and Kainz, G. (1997).
\newblock Bidirectional Heuristic Search Reconsidered.
\newblock \emph{Journal of Artificial Intelligence Research}, 7:283--317.

\bibitem[Nobes et~al.(2023)]{nobes2023voxel}
Nobes, R., Harabor, D., Wybrow, M., and Walsh, S.~D.~C. (2023).
\newblock Voxel Benchmarks for 3D Pathfinding: Sandstone, Descent, and
  Industrial Plants.
\newblock In \emph{Proc.\ Symposium on Combinatorial Search (SoCS)}.

\bibitem[Shaham et~al.(2017)]{shaham2017minimal}
Shaham, E., Felner, A., Chen, J., and Sturtevant, N.~R. (2017).
\newblock The Minimal Set of States that Must Be Expanded in a Front-to-End
  Bidirectional Search.
\newblock In \emph{Proc.\ Symposium on Combinatorial Search (SoCS)}.

\bibitem[Shaham et~al.(2018)]{shaham2018consistent}
Shaham, E., Felner, A., Sturtevant, N.~R., and Rosenschein, J.~S. (2018).
\newblock Minimizing Node Expansions in Bidirectional Search with Consistent
  Heuristics.
\newblock In \emph{Proc.\ Symposium on Combinatorial Search (SoCS)}.

\bibitem[Siag and Shperberg(2025)]{siag2025bridging}
Siag, L. and Shperberg, S.~S. (2025).
\newblock Bridging theory and practice in bidirectional heuristic search with
  front-to-end consistent heuristics.
\newblock \emph{Artificial Intelligence}, 348:104420.

\end{thebibliography}

\end{document}
